\section{Implementation}
\label{sec:impl}

\systemname\ is a Python runtime on PyTorch and HuggingFace Transformers~\cite{paszke2019pytorch,wolf2020transformers}.
Expert modules are replaced by a dispatcher that resolves identifiers to published handles.
Copies run on a dedicated CUDA stream.
A CPU thread polls completion events, publishes handles, and returns reclaimed blocks to the free list.
The fast and slow loops run on that thread.
They do not launch GPU kernels other than the copies and the expert computations the forward path already performs.
Control-period work is a scan of at most \(L\) horizon predicates plus an \(O(E_{\ell})\) score update per active layer, with \(L\) the number of MoE layers.

Both packed tiers are materialized in pinned host memory before the timed region.
Materialization proceeds one layer at a time and releases the source weights of that layer once both tiers exist, so the host does not hold the unpacked checkpoint and both packed tiers together.
A measured transition is a copy of an existing buffer.
The run aborts if any expert fails to pack, register, or fit in its pool.
There is no fallback onto unpacked weights.

The high-precision tier is FP16 and the low-precision tier is INT4 for models whose low-precision pool fits in the larger budgets we sweep.
For Qwen3-Next-80B the low tier is INT2 and the high tier is INT4, because an INT4 pool of that model already occupies most of a 48\,GB device and leaves no room for a second tier or for the KV reserve.
The tier pair is a property of the checkpoint and is held fixed across the budget sweep.
Changing \(B\) changes \(\PhiRatio\), not the quantization recipe.
Scales and zero-points travel in the handle metadata.
A layer that contains both live tiers issues two grouped GEMM launches.
The conversion workspace used by the INT4 and INT2 kernels is reserved in \(M_{\mathrm{stage}}\) as a constant number of blocks per in-flight transition, matching the workspace the kernel actually allocates on first use.
We record that size from a one-expert probe at startup rather than from a formula.

\(C\) is the ratio of copied bytes to the duration of the migration-stream event, computed only for copies that overlapped expert compute.
An idle H2D microbenchmark is recorded once per device and used solely as the alternative estimator in the model-error experiment.
\(T_{\ell}\) is the elapsed time of the expert kernels on the compute stream, which includes dequantization and both GEMM launches when both tiers are live.
The p95 miss size uses a ring of the last 64 observations of \(|\hat{U}(1)\setminus R|\) at that layer, scaled to horizon \(S\) by substituting the predicted union \(\hat{U}(S)\).
The ring is not a linear extrapolation of one layer’s count, which would double-count experts that recur across layers.

Continuous batching is supported for the latency workload.
The density \(\hat{w}\) and the miss set are properties of the batch currently in the layer, so a mixed batch is controlled by its own union of experts rather than by any single request.
Quality runs disable continuous batching and use greedy decoding so that a precision change cannot occur between the answer candidates of one multiple-choice item.
All candidates of an item are scored in one unpadded batch.

Startup classification, pool capacities, and the frozen constants \(\alpha\), \(\delta\), \(\theta_{h}\), \(T_{\mathrm{fast}}\), and \(T_{\mathrm{slow}}\) are written into the result artifact.
The artifact also records the checkpoint revision, the quantization recipe, the GPU UUID, the process peak sampled through NVML every 2\,ms, and the runtime’s own reservation counters.
A point is reportable only when those counters stayed within \(B\) for the whole run or when the log shows a refused reservation.
A run that exceeds \(B\) by allocating anyway is a bug, not a data point.

The prototype is single-GPU.
Expert-parallel placements, in which a layer’s experts are partitioned across devices, change both \(B\) and \(C\) per device and are out of scope.
The accounting identity does not depend on that restriction; the implementation does.

\paragraph{Generative-model disclosure.}
SIGMETRICS~2027 requires the methods text to state the extent of generative-model assistance.
\tbd{Authors: replace this paragraph before submission. State which prose was drafted with a language model, that the model, controller, and protocol are the authors' claims, and that every citation and measurement was checked by the authors. Do not leave this placeholder in a submission.}
